% Índice geral do curso — Meteorologia Prática
\documentclass[11pt]{article}
\usepackage{comum/preambulo}

\title{Meteorologia Prática\\
{\large Índice Geral do Curso --- baseado em R.~Stull,
\emph{Practical Meteorology} v1.02b (CC BY-NC-SA 4.0)}}
\author{}
\date{}

\begin{document}
\maketitle
\thispagestyle{fancy}

\section*{Visão geral}

Curso de graduação em física da atmosfera para estudantes com cálculo:
\textbf{22 capítulos} acompanhando o livro de Stull $+$ \textbf{1
módulo extra} (mudança climática), organizados como uma sequência de
\textbf{dois semestres} (I: Caps.~1--11; II: Caps.~12--23; ver
\texttt{guia\_do\_professor/cronograma.pdf}). Cada lei física aparece
em \emph{duas formas} — a algébrica do livro e a diferencial — e cada
capítulo fecha com um laboratório numérico em Python.

\newtcolorbox{caixaindice}{colback=corquadro!6,colframe=corquadro,
  title={Cada capítulo NN tem cinco peças},
  fonttitle=\bfseries\small,before skip=8pt,after skip=8pt}
\begin{caixaindice}
\begin{tabular}{@{}lp{0.52\linewidth}@{}}
\texttt{notas/capNN\_notas.pdf} & roteiro de quadro-negro (caixas
verdes $=$ o que escrever)\\
\texttt{slides/capNN\_slides.pdf} & apoio visual com as figuras do
livro\\
\texttt{guia\_do\_professor/guia.pdf} & seção do capítulo: ritmo,
armadilhas, demonstrações\\
\texttt{notebooks/capNN\_$\langle$tema$\rangle$.ipynb} & 4 casos
numéricos $+$ células N1--N4 do aluno\\
\texttt{notebooks/capNN\_solucoes.ipynb} & \emph{professor}: T1--T5
(com \texttt{sympy}) e N1--N4 resolvidos\\
\end{tabular}\\[3pt]
Lista impressa do capítulo: página correspondente de
\texttt{listas/listas\_exercicios.pdf}.
\end{caixaindice}

\section*{Os capítulos}

\begin{center}
\begin{tabular}{cllc}
\hline
Cap. & Título & Notebook (tema) & Sem. \\
\hline
1 & Fundamentos da Atmosfera & \texttt{fundamentos} & I \\
2 & Radiação Solar e Infravermelha & \texttt{radiacao} & I \\
3 & Termodinâmica da Atmosfera & \texttt{termodinamica} & I \\
4 & Vapor d'Água & \texttt{vapor} & I \\
5 & Estabilidade Atmosférica & \texttt{estabilidade} & I \\
6 & Nuvens & \texttt{nuvens} & I \\
7 & Processos de Precipitação & \texttt{precipitacao} & I \\
8 & Satélites e Radar & \texttt{sensoriamento} & I \\
9 & Boletins Meteorológicos e Análise de Cartas & \texttt{cartas} & I \\
10 & Forças Atmosféricas e Ventos & \texttt{ventos} & I \\
11 & Circulação Geral & \texttt{circulacao} & I \\
\hline
12 & Massas de Ar e Frentes & \texttt{frentes} & II \\
13 & Ciclones Extratropicais & \texttt{ciclones} & II \\
14 & Fundamentos de Tempestades & \texttt{tempestades} & II \\
15 & Perigos de Tempestades & \texttt{perigos} & II \\
16 & Ciclones Tropicais & \texttt{furacoes} & II \\
17 & Ventos Regionais & \texttt{ventosregionais} & II \\
18 & Camada Limite Atmosférica & \texttt{camadalimite} & II \\
19 & Dispersão de Poluentes & \texttt{poluentes} & II \\
20 & Previsão Numérica do Tempo & \texttt{previsao} & II \\
21 & Processos Climáticos Naturais & \texttt{clima} & II \\
22 & Óptica Atmosférica & \texttt{optica} & II \\
23 & Mudança Climática \emph{(módulo extra)} & \texttt{mudanca} & II \\
\hline
\end{tabular}
\end{center}

\section*{Materiais transversais}

\begin{itemize}
  \item \texttt{guia\_do\_professor/guia.pdf} --- o guia completo
        (47\,pp): uma seção por capítulo, com objetivos,
        pré-requisitos a reativar, armadilhas conceituais,
        demonstrações de baixo custo, uso dos notebooks em aula e
        pares T$\leftrightarrow$N.
  \item \texttt{guia\_do\_professor/cronograma.pdf} --- cronograma
        dos dois semestres (15 semanas cada), provas, entregas,
        avaliação sugerida e trilha compacta opcional.
  \item \texttt{listas/listas\_exercicios.pdf} --- 23 listas (uma
        página por capítulo, impressão avulsa); \emph{gerado} das
        notas por \texttt{ferramentas/gera\_listas.py}.
  \item \texttt{environment.yml} --- ambiente conda \texttt{meteo}
        (MetPy, climlab, ambiance, pyrcel, tcpyPI, Py-ART, Satpy,
        sympy).
  \item \texttt{comum/} --- preâmbulo LaTeX e tema Beamer comuns;
        \texttt{comum/estilo\_meteo.py} dá o estilo gráfico dos
        notebooks.
\end{itemize}

\section*{Para começar}

\begin{enumerate}
  \item \textbf{Professor}: leia a seção do capítulo no guia
        $\to$ prepare com as notas (as caixas verdes são o quadro)
        $\to$ projete os slides $\to$ rode os casos do notebook nos
        momentos sugeridos $\to$ libere a lista ao fechar o capítulo.
  \item \textbf{Aluno}: notas como texto-base em português $+$ o
        capítulo correspondente do Stull; notebook do capítulo para
        os casos e os exercícios N.
  \item \textbf{Ambiente}: \texttt{conda env create -f
        environment.yml}, depois \texttt{conda activate meteo} e
        \texttt{jupyter lab}.
  \item \textbf{Compilar LaTeX}: \texttt{pdflatex} (2$\times$) na
        pasta do próprio arquivo.
\end{enumerate}

\section*{Licença e atribuição}

Todo o material é derivado de R.~Stull, \emph{Practical Meteorology:
An Algebra-based Survey of Atmospheric Science} (UBC, v1.02b),
licenciado \textbf{CC BY-NC-SA 4.0} — e herda a mesma licença: uso
não comercial, com atribuição, compartilhado sob os mesmos termos.
Exceção: a Fig.~7.13 do livro (cristais de neve, © K.~Libbrecht,
``used with permission'') \emph{não} está incluída e não deve ser
redistribuída.

\end{document}
